\section{Materials}
\label{sec:materials}

In the following section the materials and solvents are characterized. 

\subsection{Solvents}

As the absorption of ZnTPP and ZnOEP is to strong to investigate them as solids, the spectroscopy is done in solution. In the following experiments, two different solvents are used, benzonitrile (bn) and toluene (tol). 

\begin{figure}[h]
    \centering
    \begin{subfigure}[b]{0.5\textwidth}
        \centering
        \includegraphics[width =0.7\textwidth]{Bilder/Grundlagen/Benzonitrile_structure.pdf}      
        \caption{Benzonitrile}
      \label{fig:BN}
    \end{subfigure}
    \hspace{\fill}
    \begin{subfigure}[b]{0.3\textwidth}
      \centering
      \includegraphics[width = 0.6\textwidth]{Bilder/Grundlagen/Toluen.pdf}      
      \caption{Toluen}
      \label{fig:Tol}
    \end{subfigure}
    \caption{Solvents used in this lab practice}
    \label{fig:Solvents}
\end{figure}

Benzonitrile is a colorless solvent with a sweet bitter almond odour. As a solvent it shows a high polarity compared to toluene due to its C -- N bound. 
Both solvents have the characteristic benzene ring. Since they are not green solvents and can cause cancer, these two solvents are replaced in industry whenever possible. 

\subsection{ZnTPP and ZnOEP}
\label{subsec:ZnTPPZnOEP}

In the following section ZnTPP and ZnOEP are introduced.

ZnTPP -- zinc tetraphenylporphyrin -- is a compound belonging to the family of porphyrins. Porphyrins are a class of 
molecules consisting of four pyrrole rings connected by methine bridges. ZnTPP contains a central zinc ion coordinated to four phenyl rings attached
to the pyrrole nitrogen atoms.
ZnTPP appears a dark purple crystalline solid. Its molecular weight is \SI{678.11}{\gram\per\mol}. It exhibits strong absorption in the visible spectrum, with a maximum absorption peak at around 558 nanometers, resulting in its characteristic purple color.
The absorption property makes ZnTPP useful in various applications, such as dye-sensitized solar cells, photovoltaic devices, and organic light-emitting
diodes. 


\begin{figure}[h]
    \centering
    \begin{subfigure}[b]{0.5\textwidth}
        \centering
        \includegraphics[width =0.7\textwidth]{Bilder/Grundlagen/ZnTPP.png}      
        \caption{5,10,15,20-Tetraphenyl-21H,23H-porphinzink(II) \newline M = 678,11 g/mol}
      \label{fig:ZnTPP}
    \end{subfigure}
    \hspace{\fill}
    \begin{subfigure}[b]{0.3\textwidth}
      \centering
      \includegraphics[width = 0.6\textwidth]{Bilder/Grundlagen/ZnOEP.png}      
      \caption{2,3,7,8,12,13,17,18-Octaethyl-21H,23H-porphinzink(II), M = 598,15 g/mol}
      \label{fig:ZnOEP}
    \end{subfigure}
    \caption{Materials used in this lab practice}
    \label{fig:Polymere}
\end{figure}

ZnOEP -- zinc octaethylporphyrin -- is also a compound belonging to the family of porphyrins. ZnOEP contains as ZnTPP a central zinc ion
coordinated to eight ethyl groups attached to the pyrrole nitrogen atoms.

One important difference for our experiments are the structural differences. As can be seen in \cref{fig:Polymere}, the benzene rings enlarge the systeme of conjugated electron system. 

